\label{sec:evaluation}

\subsection{模型的优点}
\begin{enumerate}[leftmargin=*,itemsep=0.2em]
  \item \textbf{一个模型覆盖四个问题。} 从守恒律出发建立的\eqref{eq:heat}--\eqref{eq:init} 是唯一的定解问题，四个问题只更换物性、时间范围与边界是否移动；所有结果由同一套程序、同一次云端运行产生。
  \item \textbf{离散格式具有可证明的结构性质。} 有限体积离散使水分（及常物性下的能量）平衡成为半离散系统的\textbf{精确线性不变量}（命题~\ref{prop:conserve}），并满足离散极值原理（命题~\ref{prop:extremum}）。这两条不是事后的数值观察，而是格式的性质，保证了长时间积分不会累积质量误差、不会产生非物理的负含水率。
  \item \textbf{精度可控且被独立确认。} 输出点与网格节点重合（无插值误差），收敛阶、解析解误差与 Richardson 外推共同保证四位小数有意义；校验器与求解器代码分离，二维轴对称程序是完全独立的第二实现。题目未规定之处（边界条件形式、平台外推、收缩域写法、判据分辨率、界面平均）都给出了备选结果与差异幅度（表~\ref{tab:alternatives}），使结论的稳健性可检验而不是靠断言。
\end{enumerate}

\subsection{题给模型省略的两项物理机制及其影响范围}
\label{subsec:extension}

题给参数体系把表面传质写成 $-D\partial_rC=h_m(C_s-C_{\mathrm{air}})$，即把药材的\textbf{干基含水率}与烘房空气的“水分浓度”放在同一量纲下比较，并且能量方程不含蒸发潜热。主结果必须、也确实严格按题给参数计算；本小节在题给模型之外量化这两项被省略的机制能把答案推到什么范围。全部计算在问题 3 的设定下进行（固定半径、附录 3 物性），不改变第~\ref{sec:q1}--\ref{sec:q4} 节的任何结果。

\paragraph{(a) 蒸发潜热的后验一致性诊断。} 由问题 3 的解可以反算题给传质律隐含的蒸发质量通量 $j_w=\rho_{s0}h_m(C_s-C_{\mathrm{air}})$，其中干物质表观密度 $\rho_{s0}=\rho(C_0)/(1+C_0)=\valExtDrySolidDensity$ \si{\kilo\gram\per\cubic\meter}。把其潜热负荷折算为“等效表面温降” $\Delta T_{\mathrm{lat}}=L_v(T_s)\,j_w/h$（即若要在能量上支付这一蒸发速率，表面需要比烘房低多少度）：初始时 $\Delta T_{\mathrm{lat}}=\valExtLatentDeltaTPeak$ \si{\kelvin}，\SI{3}{\hour} 时 $\valExtLatentDeltaTThreeH$ \si{\kelvin}，\SI{12}{\hour} 时 $\valExtLatentDeltaTTwelveH$ \si{\kelvin}，降到 \SI{1}{\kelvin} 以下需要 $\valExtLatentBelowOneKHours$ h（图~\ref{fig:extension}(a)）。全过程蒸发的水量为 $\valExtWaterRemovedKg$ \si{\kilo\gram} 每米长度，对应潜热 $\valExtLatentEnergyMJ$ \si{\mega\joule\per\meter}，而题给能量边界条件全程供入的净对流热量只有 $\valExtConvEnergyKJ$ \si{\kilo\joule\per\meter}（等于显热储量的增加），二者相差 $\valExtLatentToConvRatio$ 倍。\textbf{结论}：题给模型在最初十余小时显著高估药材温度（真实过程中表面接近湿球温度），因而高估 $D$、低估烘干时间；$\valExtLatentBelowOneKHours$ h 之后潜热的等效温降已不足 \SI{1}{\kelvin}，模型自洽。直接把潜热项加入能量边界条件是不可行的——题给传质律没有湿球反馈（蒸发速率不随表面降温而减小），会把表面温度压到 \SI{0}{\celsius} 以下，问题不适定。

\paragraph{(b) 能量限制（恒速干燥段）的界定分析。} 把烘房的“水分浓度”按湿空气含湿量（\si{\kilo\gram} 水/\si{\kilo\gram} 干空气）解读（饱和蒸气压用 Magnus 公式\cite{alduchov1996}，湿球温度由绝热饱和方程求根\cite{ashrae2017}），则初始状态（\SI{28}{\celsius} 与附件 1 的首个含湿量样本）对应相对湿度 $\valExtRhInitPct\%$、湿球温度 $\valExtWetBulbInit\,^\circ$C；恒温阶段（$\valQthreeAirTempPlateau\,^\circ$C、$\valQthreeAirHumPlateau$ \si{\kilo\gram\per\kilo\gram}）对应相对湿度 $\valExtRhPlateauPct\%$、湿球温度 $\valExtWetBulbPlateau\,^\circ$C（湿球温差 $\valExtWetBulbDepression$ \si{\kelvin}），数值上合理。经典干燥理论\cite{sherwood1929,mujumdar2014}指出，恒速段的蒸发通量不能超过对流供热所能维持的湿球极限
\begin{equation}
j_{\max}(t)=\frac{h\,[T_{\mathrm{air}}(t)-T_{\mathrm{wb}}(t)]}{L_v},
\label{eq:wetbulb}
\end{equation}
恒温阶段 $j_{\max}=\valExtFluxCapPlateauKgPerHour$ \si{\kilo\gram\per\square\meter\per\hour}，而题给模型的初始蒸发通量是该上限的 $\valExtFluxRatioInit$ 倍。把表面损失律改为 $\min\{h_m(C_s-C_{\mathrm{air}}),\,j_{\max}/\rho_{s0}\}$（这保持了命题~\ref{prop:conserve} 的守恒不变量，且当表面变干后自动退化为题给模型），得烘干时间 $\valExtCappedHours$ h，比主结果长 $\valExtCappedPct\%$；上限在 $\valExtCappedReleaseHours$ h 后不再起作用。把上限取 $\times0.5$ 与 $\times2$，结果分别为 $\valExtCapOneHours$ h 与 $\valExtCapThreeHours$ h（表~\ref{tab:extension}）。另一方面，若全程按湿球极限蒸发，得到纯能量下界 $\valExtEnergyBoundHours$ h（$\valExtEnergyBoundPct\%$）——主结果远高于该下界，再次说明\textbf{题给设定下的烘干时长由内部扩散控制而非能量供给控制}。

\paragraph{(c) 平衡含水率驱动力。} 物理上传质的驱动力应为 $C_s-C_{\mathrm{eq}}$，$C_{\mathrm{eq}}$ 是药材在烘房温湿度下的平衡含水率（由等温吸附曲线给出，题目未提供）。以常数 $C_{\mathrm{eq}}\in\{\valExtCeqOneValue,\valExtCeqTwoValue,\valExtCeqThreeValue\}$ 替换 $C_{\mathrm{air}}$，烘干时间分别为 $\valExtCeqOneHours$、$\valExtCeqTwoHours$、$\valExtCeqThreeHours$ h（$+\valExtCeqOnePct\%$、$+\valExtCeqTwoPct\%$、$+\valExtCeqThreePct\%$）。当 $C_{\mathrm{eq}}\to C^\ast=0.15$ 时烘干时间发散：\textbf{终点判据必须与烘房湿度相匹配}，否则物理上不可达——这是本模型给出的最有工艺价值的结论之一。值得注意的是，这一单调性并非普遍成立：驱动力减小使表面更湿，而 $D\propto\mathrm{e}^{-a/C}$ 使湿表面的扩散系数大得多，两种效应竞争；在附录 2 的律（$a=0.89$）下较大的 $C_{\mathrm{eq}}$ 反而可能略快（表面结壳减轻，case hardening），附录 3（$a=0.45$）下则单调变慢。

\input{generated/tables/tab_extension.tex}

\paperfigure{fig_extension.pdf}{模型评价扩展：(a) 题给模型蒸发率所对应的潜热等效温降 $L_vj_w/h$ 与模型实际的表面温差 $T_{\mathrm{air}}-T_s$（纵轴对数坐标，点线为 \SI{1}{\kelvin}）——前十余小时前者远大于后者，说明题给能量方程未支付蒸发潜热；(b) 题给模型、能量限制情形与三档平衡含水率驱动力下的中心水分浓度时程，点线为判据 \SI{0.15}{\kilo\gram\per\kilo\gram}。图 (a) 中模型温差在 \SIrange{2}{4}{\hour} 间因烘房温度样本的 $\pm\SI{0.3}{\celsius}$ 噪声而在零附近改变符号，负值在对数坐标下被截断，故呈现振荡带。}{fig:extension}

\subsection{局限}
\begin{enumerate}[leftmargin=*,itemsep=0.2em]
  \item \textbf{表面条件的物理内涵。} 主模型按题给参数把烘房含湿量与药材干基含水率直接相减，未使用等温吸附曲线；第~\ref{subsec:extension} 小节的 (c) 给出了这一简化的影响区间（$+\valExtCeqOnePct\%$ 至 $+\valExtCeqThreePct\%$）。
  \item \textbf{蒸发潜热未耦合。} 见 (a)、(b)；界定性结果表明真实烘干时间很可能比主结果长 $\valExtCappedPct\%$ 量级（表~\ref{tab:extension}）。注意情形 (b) 只限制质量通量而未改变能量方程（温度仍按题给模型演化），因此是\textbf{界定分析}而非完整的耦合模型。
  \item \textbf{收缩的各向同性与已知性。} 问题 4 直接使用给定的 $R(t)$，未建立含水率$\to$体积的本构关系，也未考虑轴向收缩与开裂；均匀收缩假设在出现表面硬壳时会失真。
  \item \textbf{物性与传递系数的常数化。} $h$、$h_m$ 全程恒定（收缩使表面积减小、气流条件改变时它们会变化）；$D$ 的经验公式在 $C\to0$ 时的外推缺乏数据支撑，而正是这一区间决定了烘干末段的长尾。
  \item \textbf{一维与均匀性。} 端部效应已量化为 $\valEndEffectDryDiffPct\%$（可忽略），但各向异性、非均质与批次内个体差异未建模。
  \item \textbf{烘房条件的外推。} \SI{14400}{\second} 之后的烘房状态由平台值外推（假设~\ref{as:chamber}），虽经三种取法对照（$\le\valAltPlateauLastSamplePct\%$），但若实际工艺存在升温/降湿程序，结果会改变——由弹性 $\valElasTplateau$ 可直接估算。
\end{enumerate}

\subsection{推广与工艺含义}
\begin{enumerate}[leftmargin=*,itemsep=0.2em]
  \item \textbf{模型的推广。} 本文的离散框架（Landau 坐标 $+$ 有限体积 $+$ 分段自适应 BDF $+$ 事件求根）可直接推广到：球形与板状药材（只需更换控制体测度中的几何权 $\xi^{d-1}$）；由含水率驱动的自适应收缩（把 $R(t)$ 换成由 $\bar C$ 决定的本构关系，方程形式不变）；含潜热耦合与吸附平衡的完整模型（第~\ref{subsec:extension} 小节）；以及以能耗或品质为目标的\textbf{烘房温湿度轨迹优化}——把 $T_{\mathrm{air}}(t)$ 作为控制变量、以 $t_{\mathrm{end}}$ 与热敏成分降解量为目标，本文的求解器即可作为状态方程反复调用。
  \item \textbf{工艺含义。} 灵敏度结果给出明确的优先级：(i) 提高烘房温度最有效（弹性 $\valElasTplateau$），但受热敏成分降解的约束，宜采用先低温排湿、后升温的分段程序；(ii) \textbf{减小特征尺寸最可靠}——切片或减小直径使 $t_{\mathrm{end}}\propto R^2$ 下降（问题 4 中收缩本身就把烘干时间压缩到 $1/\valQfourFixedToMovingRatio$）；(iii) 加大风速（提高 $h$、$h_m$）几乎无效（弹性 $\valElasH$、$\valElasHm$），反而增加能耗；(iv) 降低烘房湿度的收益也很小（$\valElasCair$），但必须保证 $C_{\mathrm{eq}}$ 显著低于 $C^\ast$，否则终点不可达；(v) 判据的定义值得与工艺方商榷：按“各处低于 0.15”需 $\valQthreeDryHours$ h，而此时体积平均含水率已低至 $\valQthreeMeanMoistAtDry$ \si{\kilo\gram\per\kilo\gram}，大量物料被过度干燥，能耗与品质均受损。
\end{enumerate}
